\documentclass[10pt,a4paper]{amsart}
\usepackage[margin=19mm]{geometry}
\usepackage{amsmath,amssymb,mathtools}
\usepackage[scr=rsfs,cal=euler]{mathalfa}
\usepackage{xfrac}
\usepackage[unicode,hidelinks,bookmarks=false]{hyperref}
\newcommand{\ie}{i.e.\ }
\newcommand{\cf}{cf.\ }
\newcommand{\resp}{respectively\ }
\newcommand{\Subsection}{\subsection}
\providecommand{\nobreakdash}{\nobreak\hbox{-}}
\setlength{\emergencystretch}{2em}
\multlinegap0pt
\DeclareMathOperator{\Res}{Res}
\theoremstyle{plain}
\newtheorem{theorem}{Theorem}[section]
\newtheorem{proposition}[theorem]{Proposition}
\newtheorem{lemma}[theorem]{Lemma}
\newtheorem{corollary}[theorem]{Corollary}
\newtheorem{definition}[theorem]{Definition}
\theoremstyle{remark}
\newtheorem{remark}[theorem]{Remark}
\newtheorem{example}[theorem]{Example}
\title[Evaluating $\mathrm {SU}(3)$ Verlinde sums using spectral graph theory]{Evaluating $\mathrm {SU}(3)$ Verlinde sums using spectral graph theory}
\author{Jay Jorgenson, Anders Karlsson,, Lejla Smajlovi\'c}
\date{Source: arXiv:2609.28265v1 (CC BY 4.0)}
\begin{document}
\maketitle

\noindent\emph{Source and licence.} The delivered work is the article shown in the title above, version arXiv:2609.28265v1 (CC BY 4.0), distributed under the Creative Commons Attribution 4.0 International licence, as recorded in the arXiv licence metadata for this exact version and shown on the abstract page saved with this item. The full licence notice, which carries the licence URL and the address of that abstract page, is the bundled file \texttt{licenses/arxiv-2609.28265-CC-BY-4.0.txt}.
\par\smallskip

\noindent\textit{Editorial scope.} Counted unit: the article's Theorem (source label thm:second-resultant) (source0.tex lines 1649--1666, source label \texttt{thm:second-resultant}): ``Let , which is monic of degree , and as in . Define D(w) : j 1 0 m-1 j 2 0 m-1 W(z j 1 ,z j 2 ,w). Then D(w) u (Q(u), v(Q(v),W(u,v,w)) ) and D(w) 2 m 2 w 3m-2 D(w), where is the spectral determinant p''. It is delivered together with the article's complete original proof (source0.tex lines 1668--1711, one proof environment adjacent to the statement, 44 lines), copied line for line from the article's own text. Required context retained uncounted: eq:det\_polynomial (lines 810--812); lem:constraint-product (lines 1587--1600); eq:Wdef (lines 1611--1614); lemma W (lines 1616--1625). Editorial formatting only: an AMS \texttt{amsart} preamble that restates the article's own macro definitions and its own theorem declarations, and the theorem environments the retained lines use; the article's own class file is neither shipped nor input; the retained environments are renumbered sequentially in order of appearance, whereas the article prints them with its own numbering; the article's equation numbering is not reproduced and every cross-reference inside the retained text resolves to the wrapper's numbering. No citation occurs inside any retained line, so the wrapper carries no bibliography; All retained bodies are exact copies of the bundled \texttt{sources/211567/source0.tex}, so no omission receipt applies. No asset-inclusion command occurs inside any retained span and the package ships no image asset.


\section*{Retained context: eq:det\_polynomial (source0.tex lines 810--812)}

\begin{equation}\label{eq:det_polynomial}
D(w) := \prod_{j\in G_m \setminus Z}(w-r_j).
\end{equation}


\section*{Retained context: lem:constraint-product (source0.tex lines 1587--1600)}

\begin{lemma}
\label{lem:constraint-product}
For $j_1,j_2,j_3\in\mathbb Z/m\mathbb Z$,
\[
j_1+j_2+j_3\equiv0\pmod m
\,\,\,\,\,
\text{\rm if and only if}
\,\,\,\,\,
 z_{j_1}z_{j_2}z_{j_3}=-i.
\]
Thus, for any two roots $u=z_{j_1},\,v=z_{j_2}$ of $Q(x)=x^m-i^m$,
the value $z_3:=-i/(uv)$ is automatically a root of $Q$ and is equal to
$z_{j_3}$.
\end{lemma}


\section*{Retained context: eq:Wdef (source0.tex lines 1611--1614)}

\begin{equation}
\label{eq:Wdef}
W(u,v,w) := iu^2v^2+u^2v+uv^2-2wuv+u+v-i.
\end{equation}


\section*{Retained context: lemma W (source0.tex lines 1616--1625)}

\begin{lemma}\label{lemma W}
Let $u, v$ lie on the unit circle and set $z_3 := -i/(uv)$. Then
\[
W(u,v,w) = -2uv\Bigl(w - \bigl(\operatorname{Re}(u)+\operatorname{Re}(v)+\operatorname{Re}(z_3)\bigr)\Bigr).
\]
In particular, for $u = z_{j_1}$, $v = z_{j_2}$, $j_3 = -(j_1+j_2) \bmod m$, we have that
\[
W(z_{j_1}, z_{j_2}, w) = -2z_{j_1}z_{j_2}\bigl(w - 3\Lambda_{m}(j_1,j_2,j_3)\bigr).
\]
\end{lemma}


\section{The counted unit: Theorem (source label thm:second-resultant) and its complete original proof}

\begin{theorem}
\label{thm:second-resultant}
Let $Q(u):=u^m-i^m$, which is monic of degree $m$, and $W$ as in \eqref{eq:Wdef}.
Define
\begin{equation}\label{eq:resultant_product}
\widetilde D(w) := \prod_{j_1=0}^{m-1}\prod_{j_2=0}^{m-1}W(z_{j_1},z_{j_2},w).
\end{equation}
Then
$$
\widetilde D(w) =
\Res_u\big(Q(u),\Res_v(Q(v),W(u,v,w))\big)
$$
and
\[
\widetilde D(w) = 2^{m^2}w^{3m-2}D(w),
\]
where $D(w)$ is the spectral determinant polynomial defined by \eqref{eq:det_polynomial}.
\end{theorem}

\begin{proof}
Let us write $j_3=j_3(j_1,j_2)$ viewing $j_{3}$ as a function of $j_{1}, j_{2}$ when $j_{3}:=-(j_1+j_2)\bmod m$. From Lemma \ref{lemma W} we have that
\begin{align*}
\widetilde D(w) &= \prod_{(j_1,j_2)\in\{0,...,m-1\}^2}\Big({-}2z_{j_1}z_{j_2}(w-3\Lambda_{m}(j_1,j_2,j_3))\Big)
\\&= (-2)^{m^2}\Big(\prod_{(j_1,j_2)\in\{0,...,m-1\}^2}z_{j_1}z_{j_2}\Big)\Big(\prod_{(j_1,j_2)\in\{0,...,m-1\}^2}(w-3\Lambda_{m}(j_1,j_2,j_3))\Big).
\end{align*}
There are two points to address.  First, the reduction of the product from the set of all $j_{1}, j_{2}$, and second
the evaluation of the multiplication constant.

Let $P_0:=\prod_{j=0}^{m-1}z_j$ and note that
$$
P_{0}=\prod_{j=0}^{m-1}z_j = i^m\zeta^{0+1+\dots+(m-1)}=i^m\zeta^{m(m-1)/2}.
$$
Since $\zeta^{m(m-1)/2}=e^{i\pi(m-1)}=(-1)^{m-1}$, we get
$P_0=i^m(-1)^{m-1}$. Then
$$
\prod_{(j_1,j_2)\in\{0,...,m-1\}^2}z_{j_1}z_{j_2}=P_0^{2m}=i^{2m^2}=(-1)^{m^2}.
$$
With this, we get that
$$
(-2)^{m^2}\Big(\prod_{(j_1,j_2)\in\{0,...,m-1\}^2}z_{j_1}z_{j_2}\Big)=2^{m^{2}}.
$$

As $(j_1,j_2)$ ranges the distinct $m^2$
pairs in $(\mathbb Z/m\mathbb Z)^2$,
$(j_1,j_2)\mapsto(j_1,j_2,j_3(j_1,j_2))$ is a bijection onto $A(m)$. Every pair $(j_1,j_2)$
produces a unique triple in $A(m)$, by
Lemma~\ref{lem:constraint-product}.  Therefore,
\begin{align*}
\prod_{(j_1,j_2)\in\{0,...,m-1\}^2}(w-3\Lambda_{m}(j_1,j_2,j_3)) &=
\prod_{j\in A(m)}(w-3\Lambda_{m}(j)) \\&=
\prod_{j\in A_0(m)} (w-3 \Lambda(j))\times
\prod_{j\in A(m)\setminus A_0(m)}(w-3\Lambda_{m}(j))
\end{align*}
Each $j \in A_0(m)$ is such that $\Lambda_{m}(j)=0$,
and $A_0(m)$ has $3m-2$ elements.
Therefore,
$$
\prod_{j_1,j_2}(w-3\Lambda_{m}(j_1,j_2,j_3))
= w^{3m-2}\prod_{j\in A(m)\setminus A_0(m)}(w-3\Lambda_{m}(j))
= w^{3m-2}D(w).
$$
%It remains to prove that $D$ is even polynomial. We start by noticing that every triple $(j_1,j_2,j_3)\in A(m)\setminus A_0(m)$ does not have a zero entry, hence can be uniquely paired with the triple $(m-j_1,m-j_2,m-j_3)$ which also belongs to $\in A(m)\setminus A_0(m)$. For each such pairing, using \eqref{eq. 3Lambda explicit} we have that $\Lambda_{m}(j_1,j_2,j_3)=-\Lambda_{m}(m-j_1,m-j_2,m-j_3) $. This proves that $D(w)=D(-w)$ and completes the proof.
\end{proof}

\end{document}
